\section{Theory: how a fiber field steers wave energy}\label{sec:theory}

\subsection{Plane waves in a rotated lamina}
Consider a homogeneous lamina with fiber angle $\theta$ and a plane wave
$\vu=\bm a\,e^{\ii(\omega t-k\,\vn\cdot\vx)}$ with unit wave normal
$\vn=(\cos\phi,\sin\phi)$. The engineering strain is $\veps=-\ii k\,\mathbf N(\vn)\bm a$,
with $\mathbf N(\vn)=\begin{bsmallmatrix}n_1&0\\0&n_2\\n_2&n_1\end{bsmallmatrix}$.
Substituting into \eqref{eq:strong} without damping gives the Christoffel problem
\begin{equation}
\bm\Gamma(\vn;\theta)\,\bm a=\rho c^2\bm a,\qquad
\bm\Gamma(\vn;\theta)=\mathbf N(\vn)\T\mathbf D_f(\theta)\,\mathbf N(\vn),\qquad c=\omega/k .
\label{eq:christoffel}
\end{equation}
Its two eigenvalues give the phase speeds of the quasi-longitudinal (qP) and
quasi-shear (qS) waves.

\begin{proposition}[Only the relative angle matters]\label{prop:relangle}
$\bm\Gamma(\vn;\theta)$ is unitarily similar to $\bm\Gamma(\vn_\psi;0)$, where $\vn_\psi$ is
the unit vector at the relative angle $\psi=\phi-\theta$. Hence $c=c(\psi)$, and for
an isotropic lamina $c$ is independent of $\psi$.
\end{proposition}
\begin{proof}
Let $\bm R$ be the in-plane rotation by $\theta$. Then
$\mathbf N(\bm R\vn_\psi)=\mathbf T(\theta)^{-\mathsf T}\mathbf N(\vn_\psi)\bm R^{\!\mathsf T}$
(strain rotation), and $\mathbf D_f(\theta)=\mathbf T\mathbf Q_f\mathbf T\T$. Hence
$\bm\Gamma(\vn;\theta)=\bm R\,\mathbf N(\vn_\psi)\T\mathbf Q_f\mathbf N(\vn_\psi)\bm
R\T=\bm R\,\bm\Gamma(\vn_\psi;0)\,\bm R\T$. If $\mathbf Q_f$ is isotropic, then
$\mathbf D_f(\theta)=\mathbf Q_f$ (Remark~\ref{rem:iso}) and $\bm\Gamma$ depends on $\vn$
only through its rotation.
\end{proof}

\subsection{Energy velocity and skew}
For a lossless homogeneous medium, the time-averaged energy flux of a plane wave equals
its mean energy density times the group velocity $\bm c_g=\nabla_{\bm k}\omega$
\cite{auld}. Writing $\omega=k\,c(\psi)$ in polar wave-vector coordinates,
\begin{equation}
\bm c_g=c(\psi)\,\vn+c'(\psi)\,\vn^\perp,\qquad
\tan\delta(\psi)=\frac{c'(\psi)}{c(\psi)} ,
\label{eq:skew}
\end{equation}
so the energy travels at the angle $\psi+\delta(\psi)$ from the fiber. Here $\delta$ is
the \emph{skew} between the energy direction and the wave normal. For the CFRP lamina
($E_{f1}/E_{f2}=14.6$) the qP speed falls from $9.07$ along the fiber to $2.38$ across
it, and the skew reaches $\WSkewMax^\circ$ at $\psi=75^\circ$.
\begin{itemize}[itemsep=1pt]
\item A qP wave whose normal is $60^\circ$ off the fibers carries its energy only
$4^\circ$ off them.
\item The qS skew reaches $59^\circ$.
\end{itemize}

\paragraph{Channeling.} We call the fraction of wave normals whose energy direction lies
within $20^\circ$ of the fiber the \emph{channeling share} $s_{20}$ of the lamina. For
an isotropic lamina it is $22\%$, the purely geometric fraction $20/90$. Along the
family $\mathbf Q(s)$, which keeps the isotropic invariants and scales the anisotropic
ones (Sec.~\ref{sec:dose}), it rises monotonically:
\begin{center}\small
\IfFileExists{wt_skew_table.tex}{\begin{tabular}{@{}lcccccc@{}}
\toprule
$E_{f1}/E_{f2}$ & 1.0 & 1.7 & 2.9 & 5.4 & 14.6 & 31.9\\
max qP skew & $0^\circ$ & $18^\circ$ & $35^\circ$ & $50^\circ$ & $65^\circ$ & $71^\circ$\\
$s_{20}$ & 22\% & 42\% & 60\% & 75\% & 88\% & 94\%\\
\bottomrule
\end{tabular}
}{\textit{table pending}}
\end{center}
A CFRP lamina therefore sends the energy of almost every incident direction along its
fibers. This is the physical content of the hypothesis.

\paragraph{Rays in a steered lamina, and when they apply.} When $\theta(\vx)$ varies slowly on the
wavelength scale, ray (WKB) theory applies locally. qP energy then travels along the rays of the
Hamiltonian $\omega(\vx,\bm k)=|\bm k|\,c(\psi-\theta(\vx))$, with $\psi$ the wave-normal angle and $c$
the qP phase speed of the lamina. Since $\omega$ is conserved, Hamilton's equations reduce to
\begin{equation}
\dot{\vx}=c\,\vn+c'\,\vn^\perp,\qquad \dot\psi=c'(\psi-\theta)\,\big(\vn^\perp\!\cdot\nabla\theta\big),
\label{eq:rays}
\end{equation}
the first being the group velocity \eqref{eq:skew}. Because $c'$ is small except near
$\psi-\theta\approx75^\circ$, rays are dragged along with the fibers, and the energy paths approximate
the integral curves of the fiber field. Launching rays from the drive and absorbing them in the sponge,
the target's share of the ray dwell time is a single-pass, frequency-independent estimate of its share
of the energy. Its ratio to straight fibers is a ray prediction $G_{\rm ray}$ of the redirection that
uses only the fiber field (Sec.~\ref{sec:flux}). The picture neglects the qS waves, mode conversion,
reflections and interference, and it fails where the fibers turn on the wavelength scale. These are
exactly the conditions under which we expect, and test for, departures from channeling.

\subsection{Predictions}\label{sec:predictions}
\begin{itemize}[itemsep=2pt]
\item \textbf{W1 (isotropy).} By Proposition~\ref{prop:relangle} and
Remark~\ref{rem:iso}, an isotropic lamina makes the operator independent of $\theta$.
Every orientation sensitivity vanishes identically. This is an identity of the model,
not an empirical prediction; we use it as a consistency check of the discrete gradient.
\item \textbf{W2 (redirection is the robust effect).} Write the target-energy gain of a
design against straight fibers as
\begin{equation}
G=\frac{E_f}{E_f^0}=\underbrace{\frac{E_f/E_{\mathrm{int}}}{E_f^0/E^0_{\mathrm{int}}}}_{G_\varphi\ \text{(redirection)}}
\;\times\;\underbrace{\frac{E_{\mathrm{int}}}{E^0_{\mathrm{int}}}}_{\text{field growth}},
\label{eq:gainsplit}
\end{equation}
where $E_f$ is the energy in the target and $E_{\mathrm{int}}$ the energy in the plate
interior (outside the sponge). The identity is exact. Channeling acts on
$G_\varphi$ and does not depend on frequency, because \eqref{eq:skew} is
non-dispersive. The field-growth factor depends on standing-wave interference in the
finite, lightly damped plate, and so on frequency. A gain beyond $G_\varphi$ at the
design frequency should therefore not survive a frequency change.
\item \textbf{W3 (flux follows fibers).} In any slowly curving fiber field, optimized or
not, the local energy flux should align with the local fiber director far better than
with the drive direction, and the ray prediction \eqref{eq:rays} should reproduce the redirection. The control is a set of curved fields that were never
optimized: if channeling were selected by the optimizer rather than imposed by the
lamina, the flux would not follow their fibers.
\item \textbf{W4 (anisotropy).} $G_\varphi$ should grow with the anisotropy, following
$s_{20}$, from $1$ at isotropy.
\item \textbf{W5 (limits).} A tow of minimum radius $R$ turns by at most $\ell/R$ over a run
$\ell$. Energy channeled along the fibers can therefore be redirected only as fast as
the fibers turn. Redirection should survive while tows can still turn on the
wavelength scale, and fail once $R$ exceeds it. Exact cloaking, by contrast,
needs the singular anisotropy of transformation media \cite{norris2008}. A
fixed-contrast lamina must therefore rely on interference to cancel a scattered field, and
we expect such cancellation to be tuned to one frequency (a hypothesis, tested in
Sec.~\ref{sec:cloak}, not a theorem).
\end{itemize}
